\section{What is still needed for quasi-polynomial recognition}
\label{sec:research}

The most useful next step is to make the missing global bounds explicit
enough to prove, refute, or measure. A vague expectation that every local
operation is fast is not an adequate complexity argument. In particular,
binary coordinates can encode exponentially many discs, a small grammar can
undergo exponentially many valid rewrites, and a small boundary alphabet
can carry a large algebraic object.

\subsection{A complete accounting identity}

Let $R$ be the number of processed search or hierarchy states, let $S_j$
be the complete binary encoding size of state $j$, and let $C_j$ be the
cost of generating, transforming, and checking it. If local operations have
a uniform polynomial bound $C_j\leq(S_j+N)^d$, then
\begin{equation}
 \operatorname{Time}(N)
 \leq\operatorname{preparation}(N)+
 R\left(N+\max_j S_j\right)^d.
 \label{eq:global-accounting}
\end{equation}
Thus both $R$ and the maximum represented size must be quasi-polynomial.
This elementary inequality is the central obligation that the three new
primitives do not yet discharge.

For the hierarchy route, the lexicographic reset accounting in
\cite[Section~9]{lackenby2026} gives at most $L(g+1)^L$ passes through a
step, where $L$ bounds hierarchy length and $g$ bounds relevant
pattern-complexities. Taking logarithms shows the precise requirement:
\begin{equation}
 \log L+L\log(g+1)\leq(\log(N+2))^{O(1)}
 \label{eq:reset-target}
\end{equation}
would suffice for the iteration part of a quasi-polynomial bound. A linear
bound on $L$ together with a polynomial bound on $g$ gives only an
$\exp(O(N\log N))$ reset estimate by this argument. Strict decrease of a
large integer potential proves termination, but is not the same as a small
number of decreases.

The refined hierarchy controls length through its initial handle complexity,
but this does not by itself supply \eqref{eq:reset-target}. A successful
program could either obtain much stronger simultaneous depth and complexity
bounds, or replace the lexicographic reset estimate with a better amortized
or memoized one. The latter may be more plausible than forcing every
decomposition to halve a topological quantity. In either case the actual
compressed cutting and attaching operations must satisfy the local cost
assumption in \eqref{eq:global-accounting}.

\subsection{A conditional composition theorem with explicit hypotheses}

The following statement is an accounting tool, not a theorem that the current
hierarchy or group search satisfies its hypotheses.

\begin{theorem}[Conditional quasi-polynomial composition]
Suppose a complete recognition procedure on an $N$-bit input can be organized
as a rooted search tree with the following properties, for fixed constants
$a,b,c,d,C$:
\begin{enumerate}
\item its depth is at most $C\log^a(N+2)$;
\item each node has at most $(N+2)^c$ children, including every necessary
branch of the complete procedure;
\item every node, attaching map, integer, and certificate needed locally
has total binary encoding size at most $\exp(C\log^b(N+2))$;
\item constructing a node, deciding its local actions, checking its
certificate, and verifying a leaf take time at most the $d$th power of its
encoding size plus $N$.
\end{enumerate}
Then the complete procedure has quasi-polynomial bit time and can be
executed in quasi-polynomial space by depth-first traversal.
\label{thm:conditional-qp}
\end{theorem}
\begin{proof}
The number of nodes is bounded by a geometric sum with largest term
\[
 ((N+2)^c)^{C\log^a(N+2)}
 =\exp(O(\log^{a+1}(N+2))).
\]
The per-node cost is $\exp(O(\log^b(N+2)))$ times a polynomial factor in
$N$. Multiplying gives
$\exp(O(\log^{\max(a+1,b,1)}(N+2)))$. The active path has
polylogarithmic length and each stored state has the stated bound; the local
workspace is bounded by the local running time. This gives the space
conclusion. Completeness is an explicit hypothesis; it is not
obtained from this counting argument.
\end{proof}

Another route replaces depth and branching by a proved quasi-polynomial
bound on the number of distinct canonical states and a rule that expands
each at most once. That requires a canonical representation preserving all
attaching and boundary-pattern data. Equal coarse statistics or isomorphic
unlabelled incidence graphs are not automatically equivalent states for a
future geometric continuation.

For a width-based route, $w=\log^{O(1)}N$ makes the binary tensor primitive
quasi-polynomial immediately. Planarity alone only supplies $w=O(\sqrt n)$.
Moreover, exact Jones identity still requires an independent completion
route in this implementation. A width theorem and a complete treatment of
identity cases would therefore both be needed before that observation became
a recognition theorem.

\subsection{Safe portfolio accounting}

The measurements suggest a practical improvement that already has a simple
complexity proof: run faithful Potts for a polynomially bounded amount of
whole-query work, and if it does not finish, run the certified binary tensor
algorithm from the validated diagram. The first arm is often much cheaper
on the current corpus. The second supplies the better universal bound.

\begin{proposition}[A bounded prelude preserves a fallback bound]
If an exact prelude costs at most $p(N)$ bit operations and returns either
a checked answer or an inconclusive result, while an exact complete fallback
costs at most $F(N)$, sequential composition costs at most $p(N)+F(N)$ and
is exact. In particular, a polynomially bounded faithful-Potts prelude
preserves the binary tensor full-Jones bound.
\end{proposition}
\begin{proof}
The prelude either supplies a checked answer within its budget or consumes
at most $p(N)$ before the fallback starts. The fallback receives the original
validated problem. Its cost and correctness then apply independently of the
failed prelude. Adding the bounds proves the claim.
\end{proof}

The budget must include preparation, all attempted shadings or orders,
coefficient arithmetic, and normalization, not just one inner transition
counter. In this implementation those operations have polynomial bit bounds
per represented work object, so a proved polynomial transition and state
budget can support such a prelude. The portfolio itself is proposed, not
part of the frozen patch.

There is an analogous constraint on an eventual quasi-polynomial recognizer:
it cannot always insist on finishing an uncapped $2^{O(\sqrt n)}$ invariant
first and then claim a quasi-polynomial total from a later stage. Such an
invariant must have its own quasi-polynomial allowance, be skipped, or be
proved cheaper on the relevant branch. Optional heuristics are compatible
with a good worst-case theorem when their discarded work is explicitly
bounded.

\subsection{Priority I: complete compressed geometry}

\begin{question}[Exact cutting with complete attaching data]
Can a cut along a binary normal surface be computed and independently
verified in time polynomial in the triangulation size, coordinate bit
length, and size of the compressed output, without enumerating individual
discs or parallel product regions?
\end{question}

This is the most consequential implementation target for the hierarchy
route. A useful representation must retain oriented attaching intervals,
boundary patterns, component identities, and maps back to the original
manifold. Merely recording multiplicities of local normal disc types does
not say how the new boundary pieces are paired. A benchmark should therefore
include exponentially large parallel families with a fixed number of local
types, and compare the compressed result with a literal cut on all feasible
small members. The pass condition is equality of reconstructed attaching
maps and certificate replay, together with output-size and bit-cost bounds.

\begin{question}[Certified parallelity reduction]
Can all relevant product or parallelity regions be recognized and collapsed
using compressed interval data, with a certificate of the resulting handle
structure and inherited boundary pattern?
\end{question}

This would connect the implementation to the refined hierarchy's depth
control. It must include exceptional or non-disk bundle bases and their
vertical boundaries, rather than assuming every repeated local block is a
trivial product. A concrete first milestone is an exact list of permitted
local quotient operations, each with a reconstruction certificate and
explicit effect on the chosen complexity measure. The benefit should be
measured in compressed handle count and subsequent reset work, not only in
the number of expanded tetrahedra avoided.

\begin{question}[Normal curves after a compressed cut]
Can boundary intersection, isotopy, and normalization operations be made
polynomial in all representation parameters actually produced by cutting?
\end{question}

Lackenby's fast curve algorithms provide suitable starting primitives
\cite{lackenbycurves2026}. Their hypotheses and parameters must be preserved
at the adapter boundary. In particular, standard-curve normalization depends
on its simplification number as well as triangulation size and logarithmic
weight. Controlling only binary multiplicities is therefore insufficient if
the cut creates many independently paired arc pieces. The next artifact
should expose that pairing count or simplification number as a measured and
proved state-size parameter, with literal small-instance comparison.

\begin{question}[Independent geometric decision certificates]
Can a successful hierarchy, crushing sequence, or meridian disk be exported
from discovery and replayed by a small exact checker bound to the original
diagram and its peripheral data?
\end{question}

The cocycle matching is a model for this boundary: the solver proposes a
small certificate and a separate routine checks the claimed objective.
Geometry needs richer witnesses, including the validity of the exterior,
essential boundary slopes, and the admissibility of each operation. A first
positive milestone is the Euler-characteristic disk test proved in
Section~\ref{sec:cocycle}. Its implementation should not mistake a rank-one
manifold for the input knot exterior or use a disk with inessential boundary
as an unknot certificate.

\subsection{Priority II: a global descent or reset theorem}

\begin{question}[Reset accounting stronger than lexicographic termination]
Is there an amortized potential, persistent-prefix reuse theorem, or
canonical-state bound that replaces $L(g+1)^L$ by a quasi-polynomial bound
on actual reconstruction work?
\end{question}

An informative experiment must record why a prefix is discarded, what data
remain reusable, and how much later work is repeated. Counting only successful
decompositions hides the potentially dominant cost. A candidate theorem
should specify an injection from reset events into a bounded family of
combinatorial objects or prove a multiplicative decrease after a bounded
batch of events. If neither can be shown, an explicit family with repeated
near-identical resets would be a valuable negative result.

\begin{question}[State size under repeated cutting and simplification]
Can the complete compressed state remain quasi-polynomial in the original
diagram size throughout every necessary branch?
\end{question}

The word ``complete'' includes coefficient bits, interval endpoints,
pairings, boundary-pattern labels, and certificate history. A fixed local
alphabet of handle types is helpful but does not bound the number of copies
or distinct attaching patterns. A proposed invariant should have a bound
under each primitive and under composition; an individual output-sensitive
algorithm alone does not prevent its outputs from growing too quickly.
Recording both peak size and cumulative newly created data will distinguish
reusable DAG sharing from repeated copying.

\begin{question}[Algebraic summaries beyond the Jones state space]
Can a complete detecting invariant or geometric boundary condition be
represented at polylogarithmic interfaces with quasi-polynomial total
algebraic multiplicity and update cost?
\end{question}

Replacing Bell partitions by binary states solves a specific invariant
calculation. It cannot simply be transferred to Khovanov complexes: the
number of boundary labels does not bound chain-group multiplicities,
differential data, or the work needed to perform exact cancellation and
rank decisions. Even with one boundary label, an abstract chain object can
carry arbitrarily many summands. A viable stronger summary needs an exact
composition theorem controlling those data, with proof that the final
decision is complete. Small widths on special diagram classes should be
studied separately from claims about arbitrary planar diagrams.

\subsection{Priority III: improving the new primitives}

\begin{question}[Weighted and class-aware cocycle optimization]
Can the certified flow objective be extended to local geometric costs that
predict hierarchy work, and can a bounded set of cohomology classes be
searched without losing a global bit bound?
\end{question}

Positive integral tetrahedron weights preserve the network formulation, but
require a flow algorithm polynomial in binary capacities. Minimizing genus
or pattern-complexity is not obtained simply by renaming the disc-count
objective. Varying the cohomology class also leaves the current potential
family and may reintroduce an integer search problem. A useful first result
would state exactly which weighted local objectives remain convex network
spans, and compare their certified optima with measured cut or reset costs
on genuine knot exteriors. The rank-one primitive case is a natural initial
target because its class choices are controlled.

\begin{question}[When does additional vertex freedom justify its cost?]
Can truncation or subdivision be chosen by a certified benefit/cost rule
that improves seed construction or downstream geometry on difficult inputs?
\end{question}

The finite-exterior examples show useful freedom with ten vertices and none
after one-vertex simplification. More vertices do not automatically help:
they increase the flow and geometric state, and disc count changes with the
triangulation. A fair comparison should report the number of tetrahedra,
bits, normal discs, Euler characteristic, and actual downstream work. For
known pure gauge changes on a fixed triangulation, the certificate-transport
proposition avoids solving the flow again. Retriangulation requires a
separate transport theorem and cannot use that formula without proof.

\begin{question}[Smaller tensor work with the same certified bound]
Can order selection, outer-face selection, and a bounded faithful prelude
reduce actual state and arithmetic costs while preserving the universal
square-root bound?
\end{question}

The existing certificate allows heuristic orders to compete with a guaranteed
one. A cost model incorporating actual state multiplicity and mantissa bits
may predict runtime better than frontier width alone. Any order-search
budget must be charged. An especially useful negative deduction follows
from Proposition~\ref{prop:gauge}: changing only a same-outer bend gauge
cannot reduce normalized mantissas or their bit lengths for a fixed order.
Optimizing that gauge solely to shrink integer precision would be wasted
effort. Changing the outer face is a different operation and can be tested
explicitly, while retaining the complete turning proof. Sharper bounds on
formal coefficient span also remain a mathematical question, rather than
an assumed consequence of gauge covariance.

\begin{question}[Witness quality and the next compressed bottleneck]
Can a budgeted overlap scheduler obtain useful gain guarantees, while
rejecting impossible whole-donor queries before they consume the allowance?
\end{question}

The present first-witness gain ratio tends to zero on a proved family, so
an unconditional approximation theorem for this exact policy is unavailable.
A modified policy might retain several cheap witnesses, increase a threshold
geometrically, or optimize gain only within a controlled collection of
periodic alignments. Each proposal needs a clear extra-work bound and a
counterexample search. Independently, the pair $(ab)^N,(aabb)^N$ calls for
a cheap exact incompatibility summary of run lengths or periodic defects
at the earlier whole-donor stage. The relevant parameter may be defect count
in addition to grammar size and length bit size.

\begin{question}[From local group moves to complete bounded discovery]
For presentations arising from actual knot diagrams, can one prove a
quasi-polynomial bound on necessary moves, grammar growth, and certificate
size for a complete discovery strategy?
\end{question}

This is substantially harder than a polynomial local word theorem. A strict
decrease of expanded relator length may still allow exponentially many
unit decreases. Strong compression does not itself make the greedy move
set complete. Research should distinguish arbitrary synthetic presentations
from residuals reachable by verified transformations of Wirtinger
presentations. The next benchmark target is a collection of actual knot
residuals above the explicit bridge cap that genuinely exercise the new
partial-overlap stage, with independently replayed final certificates and
all preceding-stage work included.

\subsection{Milestones and decision criteria}

The priorities above can be organized into three reviewable releases.
The first should implement a capped full-Jones portfolio, seed-certificate
transport, and a peripheral-data-bound positive disk checker. Success means
unchanged exact outputs, independent certificate replay, and a documented
improvement on inputs that exercise the changed path.

The second should deliver compressed cutting and normal-curve adapters with
explicit output-size bounds, including product-region and boundary pairing
data. Success means agreement with literal topology on small cases and
polynomial dependence on binary weights on large structured families, with
no hidden disc expansion. If some family forces excessive pairing growth,
that obstruction should be stated before extending the hierarchy code.

The third should address \eqref{eq:reset-target} and
\eqref{eq:global-accounting}: prove an amortized reset or canonical-state
bound for a complete strategy, or construct a rigorous family showing that
the proposed strategy cannot satisfy it. A general quasi-polynomial claim
should be made only when the completeness theorem, state-size invariant,
primitive bit bounds, and total-work accounting all refer to the same
implemented or fully specified algorithm.

The work in this package makes the next stage concrete. It provides exact
primitives whose guarantees can be composed when the missing global
structure is available, and it identifies several attractive shortcuts
that fail under direct mathematical or experimental scrutiny.
